\section{Method}
\label{sec:method}

Our design philosophy is \emph{intervention, not competition}: rather than
proposing a new system, we hold a strong, widely used SELD backbone fixed. We
then vary three binary factors in a fully crossed factorial design: the geometry
prior (the variable of interest), the temporal architecture (how the network
integrates the multichannel array signal over time), and the input modality. This
lets any metric change be attributed to a named factor rather than to a tangle of
co-varying design choices. This section describes the shared backbone
(Sec.~\ref{sec:backbone}), the one-bit geometry intervention
(Sec.~\ref{sec:gca}), the temporal backbones (Sec.~\ref{sec:arch}), and the
metrics, statistics, and probing protocol (Sec.~\ref{sec:stats}).

\subsection{Shared backbone}
\label{sec:backbone}
We adopt the official DCASE~2024 Task~3 baseline~\cite{krause2024sde} as the
fixed scaffold. Audio sampled at $24$\,kHz is segmented into $5$-second
($250$-frame) inputs. The two input modalities differ only in their channel
composition: the four-channel \emph{microphone-array} (MIC) input is rendered as
ten feature maps ($4$ log-mel spectrograms $+$ $6$ GCC-PHAT cross-correlation
maps), while the \emph{first-order Ambisonics} (FOA) input is rendered as seven
maps ($4$ log-mel spectrograms $+$ $3$ acoustic intensity-vector maps). Both
modalities thus carry their spatial information through standard, implicitly
geometric features; the explicit prior we ablate is added on top.

A shared three-block 2-D convolutional stack ($64$ filters per block, batch
normalization, ReLU) with frequency pooling factors $[4,4,2]$ and a single
temporal pooling of $5$ reduces each input to a $50$-frame sequence of
$128$-dimensional embeddings. A temporal model (Sec.~\ref{sec:arch}) then mixes
this sequence over time before per-frame Multi-ACCDDOA
heads~\cite{shimada2022multiaccdoa} emit $3$ tracks $\times$ $13$ classes
$\times$ (activity, $3$-axis Cartesian DOA, distance), trained with the auxiliary
duplicating permutation-invariant (ADPIT) loss. All cells start from the official
synthetic-data-pretrained weights and fine-tune for $60$ epochs with the
\emph{identical} optimizer (Adam, learning rate $10^{-3}$, batch size $32$,
dropout $0.05$); we verify baseline parity against the published numbers in
Section~\ref{sec:repro}. Because the convolutional front-end, heads, loss,
optimizer, data, and width are held fixed across every cell, capacity and recipe
are controlled and the only differences are the three factors we vary.

\subsection{The geometry prior as a one-bit intervention}
\label{sec:gca}
The geometry-aware channel-attention (GCA) block performs single-head attention
\emph{over the microphone-channel dimension} of the input feature tensor, i.e.\
upstream of the shared convolutional front-end, letting the network reweight and
recombine the array channels before convolution. Operating on the raw array
channels rather than post-convolution feature maps is what makes a
microphone-pair geometry bias well-defined. Its geometry-aware variant injects a
fixed, layout-derived feature $g_{ij}\in\mathbb{R}^4$ for each ordered channel
pair $(i,j)$ into the attention \emph{keys}, so that geometry shapes the score
before the dot product rather than as a constant added afterwards:
\begin{equation}
  \mathrm{att}(i,j) = \frac{q_i^\top\!\bigl(k_j + W_g\,g_{ij}\bigr)}{\sqrt{d}}
  = \frac{q_i^\top k_j}{\sqrt{d}} + \frac{q_i^\top W_g\,g_{ij}}{\sqrt{d}},
  \label{eq:gca}
\end{equation}
where $q_i,k_j\in\mathbb{R}^d$ are the learned query/key projections of the
per-channel embeddings, $d$ is their dimension, and $W_g\in\mathbb{R}^{d\times4}$
(with bias) is the only learned geometry parameter. The pair feature $g_{ij}$ is a
deterministic, non-trainable function of the array layout (a registered buffer):
for the microphone array it is $(\Delta x,\Delta y,\,\text{dist},\,\text{bearing})$
between mics $i$ and $j$ in the array plane, and for Ambisonics it is the
difference of the two channels' direction-of-maximum-response vectors together with
a both-directional indicator, so the omnidirectional $W$ channel and the $X/Y/Z$
first-order dipoles are correctly distinguished. Because the geometry enters only
through the keys, its contribution to the logit, $q_i^\top W_g\,g_{ij}/\sqrt{d}$,
is \emph{content-dependent} (it is modulated by the query) rather than a fixed
additive bias. The prior therefore tells the channel-attention
layer which channels are spatially related \emph{before} it has seen any data.

This block is the locus of our single intervention. The one-bit toggle is our
sole design variable: \textbf{\texttt{full}} keeps the geometry projection $W_g$
active, so Eq.~\eqref{eq:gca} is geometry-aware; \textbf{\texttt{no\_geom}} drops
the $W_g\,g_{ij}$ term, reducing Eq.~\eqref{eq:gca} to plain content-based channel
attention. Crucially, the two variants differ by exactly the $80$ parameters of
$W_g$ (a $\mathrm{Linear}(4,16)$ projection: $4\times16$ weights plus $16$ biases)
and are otherwise bit-for-bit identical in architecture, initialization, data,
and optimization, so a paired \texttt{full}$-$\texttt{no\_geom} contrast isolates
the \emph{geometric prior} with essentially no change in capacity. To further
separate the geometric bias from the attention \emph{mechanism} that carries it,
each cell also includes a \textbf{no-GCA} control that removes channel attention
entirely, and the MIC+CRNN cell additionally includes a vanilla
squeeze-and-excitation control~\cite{hu2018senet}; both confirm that it is the
geometry, not channel attention per se, that produces the effects we report. The
toggle is well-defined and applied identically in both modalities, making the
prior directly comparable across MIC and FOA.

\subsection{Three temporal backbones spanning the locality spectrum}
\label{sec:arch}
The hypothesis we test is that the geometry prior's usefulness depends on how
much \emph{temporal locality} the backbone already assumes
(Fig.~\ref{fig:locality}). We therefore swap \emph{only} the temporal model that
sits between the shared convolutional front-end and the shared heads, choosing
three backbones that span the locality--globality spectrum while being matched in
width ($128$) and depth ($\approx\!4$ temporal layers):

\begin{figure*}[t]
\centering
\includegraphics[width=0.92\textwidth]{locality_spectrum.png}
\caption{The temporal-locality spectrum and our hypothesis. As the temporal
backbone changes from a recurrent CRNN (strong built-in locality) through a
convolution--attention Conformer (hybrid) to a pure-attention Transformer
(global all-to-all mixing), the value of the hand-crafted array-geometry prior is
predicted---and measured---to move monotonically from helpful, through
near-neutral, to harmful for directional accuracy. Numbers are the in-domain extreme effects on
DOAE$_\mathrm{CD}$ (Sec.~\ref{sec:main}).}
\label{fig:locality}
\end{figure*}
\begin{itemize}
  \item \textbf{CRNN+MHSA} (the official baseline), the most \emph{local}
  backbone: a two-layer bidirectional GRU~\cite{cho2014gru} whose recurrence
  imposes a strong recency/locality bias and limits arbitrary channel
  recombination, followed by two layers of eight-head
  self-attention~\cite{sudarsanam2021attention}.
  \item \textbf{Conformer}, the \emph{hybrid}: a linear projection followed by
  four Conformer blocks~\cite{gulati2020conformer}, each interleaving eight-head
  self-attention with a depthwise-convolution module (kernel $31$) and
  feed-forward layers (dimension $512$). It restores locality \emph{within} an
  attention model, and so is expected to behave between the two extremes.
  \item \textbf{Transformer}, the most \emph{global} backbone: a linear
  projection followed by a four-layer pre-norm Transformer
  encoder~\cite{vaswani2017attention} with eight heads, which removes recurrence
  and convolution-in-time entirely and mixes all frames and channels globally
  through self-attention.
\end{itemize}
All three consume the identical $50$-frame convolutional sequence and feed the
identical Multi-ACCDDOA heads, so the comparison isolates temporal mixing. The
ordering CRNN $\to$ Conformer $\to$ Transformer is thus a monotonic decrease in
built-in locality, which is the axis along which we predict the prior's effect to
vary. Crossing the three architectures with the two modalities gives six cells;
within each, the geometry toggle and its controls yield the task identifiers in
Table~\ref{tab:grid}. Every cell is run with five random seeds.

\begin{table}[t]
\centering
\caption{The grid and task IDs. ``full''/``no\_geom'' is the geometry-bias
toggle; controls remove channel attention. All cells use five seeds.}
\label{tab:grid}
\begin{tabular}{llccc}
\toprule
Architecture & Modality & \texttt{full} & \texttt{no\_geom} & control \\
\midrule
CRNN+MHSA   & MIC & 110 & 111 & 112 / 113 \\
CRNN+MHSA   & FOA & 130 & 131 & 100 \\
Conformer   & MIC & 161 & 162 & 160 \\
Conformer   & FOA & 171 & 172 & 170 \\
Transformer & MIC & 141 & 142 & 140 \\
Transformer & FOA & 151 & 152 & 150 \\
\bottomrule
\end{tabular}
\end{table}

\subsection{Metrics, statistical protocol, and probing}
\label{sec:stats}
\textbf{Metrics.} We report the official DCASE~2024 metrics: location-aware
$F_{20\degsym}$ (detection $F$-score counting a detection correct only within
$20\degsym$), class-dependent localization error DOAE$_\mathrm{CD}$ (degrees),
location recall, relative distance error (RDE), and the aggregate SELD score. Our
headline metric is DOAE$_\mathrm{CD}$---the directional accuracy that is the core
figure of merit for any direction-of-arrival estimator and exactly what an
array-geometry prior is meant to improve; RDE serves as the orthogonal distance
axis used to test whether the architecture that modulates direction also
modulates distance.

\textbf{Paired contrasts.} Within each cell we match the \texttt{full} and
\texttt{no\_geom} runs by seed and summarize the geometry effect as the paired
difference $\Delta=\text{full}-\text{no\_geom}$, reported with a paired $t$-test,
a Wilcoxon signed-rank test, the standardized effect size Cohen's \dz, and a
bias-corrected bootstrap $95\%$ confidence interval ($5000$ resamples).

\textbf{Primary test.} Because individual per-cell contrasts have limited power
at $n=5$, our hypothesis does not rest on any single $t$-test, but accounts for
variance explicitly rather than over-interpreting a few
runs~\cite{bouthillier2021variance}. Our \emph{primary} test is a factorial
Type-II ANOVA over the per-seed metrics that asks whether the geometry effect
\emph{interacts} with the temporal architecture. The focal test contrasts the two
locality extremes (recurrent vs.\ pure-attention) in a
modality$\times$architecture$\times$prior model. Because the same five seeds are
matched across cells, we corroborate this interaction with a within-seed paired
second-difference test (and confirm with a per-seed random-intercept mixed model
that no extra between-seed variance is present), and we test the full graded
hypothesis with an ordered Jonckheere--Terpstra trend test across all three
backbones rather than only the extremes. We also fit the unordered three-level
factorial including the Conformer, and report all of these. We do not correct
across the secondary per-cell contrasts, which we treat as descriptive; the
confirmatory claim is the pre-specified architecture$\times$prior interaction and
its monotone trend.

\textbf{Power and the $n=5$ budget.} Five seeds is a deliberate compute
trade-off---the grid spans six cells plus no-GCA controls, each retrained and
evaluated on two external corpora---and the paired design recovers most of the
power a small $n$ would cost, because the large between-seed variance cancels in
the within-seed difference $\Delta$. Our confirmatory statistics are therefore
defined at the \emph{grid} level: the focal interaction uses all $40$ per-seed
observations of the recurrent/Transformer sub-grid and the trend test all $60$,
targeting the large effects we observe ($|\dz|\approx0.7$--$1.0$ at the extremes,
for which a paired $n=5$ test retains usable power) rather than small per-cell
differences. We treat the independent-dataset replication (Section~\ref{sec:cross})
as the decisive evidence rather than any single in-domain $p$-value.

\textbf{External validity.} Every effect is re-examined out of distribution. We
evaluate the STARSS23-trained checkpoints zero-shot on the independent STARSS22
recordings~\cite{politis2022starss22} (distance-blind, as STARSS22 lacks distance
labels) and stress-test them on the synthetic TAU-NIGENS-SSE-2021
corpus~\cite{politis2021nigens}, mapping the overlapping classes to the STARSS
taxonomy. An effect that reproduces on an independent recording campaign is far
stronger evidence than any in-domain $p$-value.

\textbf{Mechanism probe.} To ask \emph{why} a cell behaves as it does, we linearly
probe the representation entering the temporal model. Post-convolutional features
are pooled over frequency, and a ridge regressor predicts $(\sin,\cos)$ of
azimuth and elevation on single-source frames using five-fold cross-validation
split by recording. The resulting angular MAE measures how \emph{linearly
decodable} direction is in the shared representation; comparing \texttt{full}
against \texttt{no\_geom} isolates the prior's effect on linear representation
content, independent of the decoder, and tests whether any behavioural effect is
mirrored by a change in linearly accessible spatial information.
